\documentclass[12pt,a4paper]{article}
\usepackage[utf8]{inputenc}
\usepackage[spanish,es-noquoting]{babel}
\usepackage[T1]{fontenc}
\usepackage{lmodern}
\usepackage{geometry}
\geometry{top=2cm, bottom=2cm, left=2.5cm, right=2.5cm}
\usepackage{xcolor}
\usepackage{listings}
\usepackage{enumitem}
\usepackage{fancyhdr}
\usepackage{hyperref}
\usepackage{tikz}
\usetikzlibrary{arrows.meta, positioning, calc, shapes.geometric}

\definecolor{codegreen}{rgb}{0,0.6,0}
\definecolor{backcolour}{rgb}{0.95,0.95,0.92}

\lstset{
    language=C,
    backgroundcolor=\color{backcolour},
    commentstyle=\color{codegreen},
    keywordstyle=\color{blue},
    basicstyle=\ttfamily\small,
    breaklines=true,
    frame=single,
    numbers=left,
    tabsize=4,
    extendedchars=true,
    showstringspaces=false,
    inputencoding=utf8,
    literate={á}{{\'a}}1 {é}{{\'e}}1 {í}{{\'i}}1 {ó}{{\'o}}1 {ú}{{\'u}}1 {ñ}{{\~n}}1 {¡}{{\textexclamdown}}1 {¿}{{\textquestiondown}}1
}

\pagestyle{fancy}
\fancyhf{}
\lhead{Monitorías Sistemas Operativos 2026-I}
\rhead{Capítulo 3: Creación de Hilos}
\cfoot{\thepage}

\title{\textbf{Guía de Aprendizaje: Introducción a los Hilos (POSIX Threads)}}
\author{Malak Sanchez \\ \small Monitorías de Sistemas Operativos}
\date{\today}

\begin{document}

\maketitle

\section{Motivación}
Hasta ahora, hemos visto procesos que se ejecutan de forma independiente. Sin embargo, crear un proceso nuevo con \texttt{fork()} es costoso porque el sistema operativo debe copiar casi todo el contexto del padre. Los \textbf{hilos} (threads) son unidades de ejecución más ligeras que conviven dentro de un mismo proceso y comparten gran parte de sus recursos.

\section{Idea Fundamental}
Un hilo es el ``flujo de ejecución''  más pequeño. Un solo proceso puede tener múltiples hilos ejecutándose simultáneamente. La diferencia clave es:
\begin{itemize}
    \item \textbf{Procesos:} Tienen su propio espacio de memoria (aislamiento).
    \item \textbf{Hilos:} Comparten el código, los datos globales y el heap, pero tienen su propia pila (stack) y contador de programa (PC).
\end{itemize}

\begin{center}
\begin{tikzpicture}[font=\small, >=Stealth,
    container/.style={draw, thick, fill=gray!5, minimum width=6cm, minimum height=4cm, align=center},
    thread/.style={draw, thick, fill=blue!10, rounded corners=3pt, minimum width=1.5cm, minimum height=2cm, align=center}]

    % --- PROCESS CONTAINER ---
    \node[container] (proc) at (0,0) {};
    \node[above] at (proc.north) {\textbf{Proceso (Recursos Compartidos)}};
    
    % Shared segments
    \node[draw, fill=yellow!20, minimum width=5.5cm] at (0, 1.3) {Code / Data / Heap / Files};

    % Threads
    \node[thread] (t1) at (-1.5, -0.5) {Thread 1 \\ \tiny (Stack 1)};
    \node[thread] (t2) at (1.5, -0.5) {Thread 2 \\ \tiny (Stack 2)};

    % Execution wave
    \draw[blue!60!black, ultra thick, out=45, in=225] (-1.8, -1) to (-1.2, 0);
    \draw[blue!60!black, ultra thick, out=45, in=225] (1.2, -1) to (1.8, 0);

\end{tikzpicture}
\end{center}

\newpage

\section{Sintaxis Básica (pthreads)}
Para trabajar con hilos en Linux usamos la librería \texttt{pthread} (hay que compilar con \texttt{-lpthread}).

\begin{lstlisting}[numbers=none]
#include <pthread.h>

// 1. Definition of the function the thread will run
void* thread_routine(void* arg);

// 2. Variable to store the thread ID
pthread_t thread_id;

// 3. Thread creation
pthread_create(&thread_id, NULL, thread_routine, (void*)arg);
\end{lstlisting}

\section{Ejemplo Mínimo: Hola desde un Hilo}
\begin{lstlisting}
#include <stdio.h>
#include <pthread.h>
#include <unistd.h>

void *say_hello(void* arg);

int main(){
    pthread_t thread;

    // Create the thread
    if(pthread_create(&thread, NULL, say_hello, NULL) != 0){
        perror("Failed to create thread");
        return 1;
    }

    printf("Main thread: Thread created successfully.\n");
    
    // Wait for the thread (we will deepen this in section 2)
    pthread_join(thread, NULL); 

    return 0;
}

void *say_hello(void* arg){
    printf("Greetings from the new thread!\n");
    return NULL;
}

\end{lstlisting}

\section{Compilación}
Es fundamental añadir el flag \texttt{-lpthread} al compilar con \texttt{gcc}:
\begin{lstlisting}[language=bash, numbers=none]
gcc main.c -o program -lpthread
\end{lstlisting}

\section{Ejercicios}
\begin{enumerate}
    \item Escriba un programa que cree dos hilos. Cada hilo debe imprimir un mensaje diferente 5 veces.
    \item ¿Qué sucede si el hilo principal (\texttt{main}) termina antes de que los otros hilos finalicen? Investíguelo.
\end{enumerate}

\end{document}
